\section{Density Policy Design}
\label{app:density_policy_design}

This section details the sampling policies used in
Tables~\ref{tab:standard_benchmarks} and~\ref{tab:industrial_benchmarks}.
Although the refinement signals are task dependent, all policies follow the
same pipeline: construct a refinement score, convert it into a target sampling
density, sample the prescribed number of points, and estimate the effective
density represented by the selected samples for measure correction.

\paragraph{Target sampling density.}
Let
\[
    \rho=\frac{1}{N}\sum_{i=1}^{N}\delta_{x_i}
\]
denote the reference point measure and let $n$ be the inference sampling
budget. Each policy produces a nonnegative refinement score $r_\psi(x)$ and
a corresponding positive sampling weight $t_\psi(x)$. We normalize this weight
with respect to the reference measure,
\begin{equation}
    \lambda_{\mathrm{tar}}(x)
    :=
    \frac{t_\psi(x)}
         {\int_\Omega t_\psi(y)\,\mathrm{d}\rho(y)},
    \qquad
    \mathrm{d}\widetilde{\rho}_{\mathrm{tar}}(x)
    =
    \lambda_{\mathrm{tar}}(x)\,\mathrm{d}\rho(x).
    \label{eq:app_target_density}
\end{equation}
Thus, $\lambda_{\mathrm{tar}}$ directly specifies where the inference samples
should be concentrated. The additive construction
$t_\psi=b_0+a_0r_\psi$ is equivalent, after normalization, to the mixture
formulation in \eqref{eq:unified_refinement_measure}.

\paragraph{Coverage-based measure correction.}
The target density specifies how points are sampled, while the selected samples
may represent unequal portions of the reference discretization. We therefore
estimate the effective coverage of each retained point and use its inverse
density for the measure correction in Physics-Attention. If a selected point
$x_i$ represents $m_i$ reference points, its correction weight is
\begin{equation}
    c_i
    =
    \frac{m_i}{N/n}.
    \label{eq:app_coverage_correction}
\end{equation}
For region-based sampling, the same principle gives
$c_i\propto N_k/n_k$, where $N_k$ and $n_k$ are the reference and retained
point counts in region $P_k$. A common scale in $c_i$ cancels in the
self-normalized aggregation. 

\paragraph{Shared refinement features.}
The three policies differ primarily in how they construct $r_\psi$.
We use local geometric or physical priors for PI, variation in a coarse
prediction for CPG, and learned intermediate features for CA.
For solution-dependent policies, local variation is estimated using
neighbor differences, local-mean residuals, or least-squares gradients.
When the score is available only on a coarse set, it is transferred to the
reference discretization by nearest-neighbor, inverse-distance, or linear
interpolation. Surface signals can additionally be extended to nearby volume
points with distance-decaying weights.  For compactness, we summarize the shared refinement and transfer operators in Table~\ref{tab:density_primitives}. These primitives are combined differently across tasks, while all policies ultimately construct a target density \(\lambda_{\mathrm{tar}}\).

\paragraph{Sampling procedures.}
Depending on the geometry, the target density is realized either by direct
weighted sampling or by regional sampling. The latter partitions the domain
into spatial regions, distributes the sampling budget according to regional
scores, and places points within each region using farthest-point sampling.
For CPG and CA, the coarse samples $S_0$ are retained; only the remaining
budget is sampled from the residual target density after accounting for the
coverage already provided by $S_0$. The following subsections specify the
task-specific refinement scores and density constructions.

\begin{table}[t]
\centering
\caption{Shared primitives used to construct the sampling densities.}
\label{tab:density_primitives}
\small
\begin{tabular}{@{}p{0.13\linewidth} p{0.54\linewidth} p{0.25\linewidth}@{}}
\toprule
Notation & Definition & Role \\
\midrule

$D_k[f]_i$
&
$\displaystyle
\frac{1}{k}\sum_{j\in\mathcal N_k(i)}
\frac{\|f_i-f_j\|_2^2}
     {\|z_i-z_j\|_2^2+\epsilon}
$
&
Local gradient-energy score
\\[1.2em]

$L_k[f]_i$
&
$\displaystyle
\left\|
f_i-\frac{1}{k}\sum_{j\in\mathcal N_k(i)}f_j
\right\|_2
$
&
Local curvature / graph-Laplacian proxy
\\[1.2em]

$G_k[f]$
&
$\displaystyle
\|\nabla_{\mathrm{LSQ}}f\|_F
$
&
Least-squares gradient magnitude
\\[0.8em]

$\mathcal Q_q[f]$
&
$\displaystyle
\frac{f}{Q_q(f)}
$
&
Quantile-based normalization
\\[1em]

$\mathcal T_{a,b}[f]$
&
$\displaystyle
\operatorname{clip}_{[0,1]}
\left(
\frac{f-Q_a(f)}
     {Q_b(f)-Q_a(f)}
\right)
$
&
Robust rescaling and clipping
\\[1.2em]

$\mathcal I_0,\mathcal I_8,\mathcal I$
&
Nearest-neighbor, eight-neighbor inverse-distance,
and linear interpolation, respectively
&
Transfer coarse scores to reference points
\\[0.8em]

$\mathcal E_\ell[v](x)$
&
$\displaystyle
e^{-d_\Gamma(x)/\ell}v(x_\Gamma)
$
&
Extend a surface score to nearby volume points
\\[1em]

$\mathcal K_\ell[v](x)$
&
$\displaystyle
\max_{i\in\Gamma}
e^{-\|x-x_i\|/\ell}\,
\mathcal Q_{0.95}[v]_i
$
&
Propagate localized surface importance
\\[1.2em]

$\mathcal P_{K,\kappa}[s](x)$
&
$\displaystyle
\frac{n_k}{n\mu_k},
\qquad x\in P_k
$
&
Region-based sampling density
\\

\bottomrule
\end{tabular}
\end{table}

\subsection{Physics-Informed Density Refinement}
\label{app:pi_density_design}

PI constructs densities from geometry or input coefficients. For Pipe and ShapeNet-Car, we first sample a fraction $1-\gamma$ of the budget uniformly across the reference mesh, then use the remaining fraction $\gamma$ near the wall. Let $H_d$ contain points within wall distance $d$, so $\rho(H_d)=|H_d|/N$ is their fraction of the reference mesh. The target density combines global coverage with near-wall refinement:
\begin{equation}
    \lambda_{\mathrm{tar}}(x)
    =\underbrace{(1-\gamma)}_{\text{global coverage}}
     +\underbrace{\gamma\frac{\mathbf{1}_{H_d}(x)}{\rho(H_d)}}_{\text{near-wall refinement}}.
    \label{eq:app_density_shell}
\end{equation}
Thus, points near the wall receive additional density while the rest of the mesh retains a uniform base. Budget fractions are rounded to integer point counts. If the near-wall candidates are exhausted, we fill the remaining budget with the closest available points outside the wall band.

\paragraph{Elasticity.}
The internal hole induces stress concentrations, motivating refinement near its boundary. We estimate the cavity distance $d_{\mathrm{cav}}$ from empty regions of a $64\times64$ auxiliary grid and let the additional density decay exponentially with distance:
\begin{equation}
    \lambda_{\mathrm{tar}}=\mathcal{N}[1+64\exp(-d_{\mathrm{cav}}/0.012)].
    \label{eq:app_pi_elasticity_density}
\end{equation}

\paragraph{Plasticity.}
Die contact and changes in the die profile localize deformation, so we refine near the contact region and beneath curved parts of the die. Let $z\in[0,1]$ denote normalized depth from contact and $k_d=|\partial_{xx}S_d|$ the profile's second-difference magnitude, with profile mean $\overline{k_d}$. We combine depth decay with this curvature indicator:
\begin{equation}
    \lambda_{\mathrm{tar}}
    =0.2+0.8\mathcal{N}\!\left[e^{-z/0.1}\left(1+0.5\frac{k_d}{\overline{k_d}}\right)\right].
    \label{eq:app_pi_plasticity_density}
\end{equation}

\paragraph{Airfoil.}
We use distance to the airfoil surface as a prior because surface pressure determines aerodynamic forces. All 121 surface nodes are retained, and the remaining points follow an exponential distance weight. With $d_\Gamma$ denoting physical distance to the surface, both metrics use
\begin{equation}
    \lambda_{\mathrm{tar}}^{\mathrm{field}}
    =\lambda_{\mathrm{tar}}^{C_D/C_L}
    =\mathcal{N}[1+20e^{-d_\Gamma/0.3}].
    \label{eq:app_pi_airfoil_density}
\end{equation}

For $C_D/C_L$, we screen 127 combinations of prior amplitude, distance scale, and density compensation on 24 calibration geometries with one sampling seed, then select among the shortlisted settings on 46 geometries with three seeds. The selected rule uses index-grid Voronoi compensation with $\beta=1$ and matches the existing field PI policy. One setting is used across all three budgets, followed by evaluation on the separate 48-geometry confirmation pool and the 110 test geometries.

\paragraph{Pipe.}
The no-slip walls produce near-wall velocity gradients, motivating additional resolution there. We reserve $75\%$ of the budget for uniform coverage and $25\%$ for points within wall distance $0.05$:
\begin{equation}
    \lambda_{\mathrm{tar}}(x)=0.75+0.25\frac{\mathbf{1}_{H_{0.05}}(x)}{\rho(H_{0.05})}.
    \label{eq:app_pi_pipe_density}
\end{equation}

\paragraph{Darcy.}
Permeability jumps mark changes in the pressure gradient, motivating refinement near coefficient interfaces. Let $\mathcal{I}_a$ contain nodes adjacent to a coefficient jump and define the implemented distance proxy $d_a(x_i)=\min_{j\in\mathcal{I}_a}|i-j|/84$ along the flattened reference-grid ordering. The density decays away from these interfaces:
\begin{equation}
    \lambda_{\mathrm{tar}}=\mathcal{N}[1+0.25\exp(-d_a/(2/84))].
    \label{eq:app_pi_darcy_density}
\end{equation}

\paragraph{ShapeNet-Car.}
Surface pressure and near-wall velocity gradients determine drag, so the prior favors the vehicle surface and nearby fluid. We augment a uniform base with points from $H_{0.1}$. This near-wall sampling stage uses $50\%$ of the budget for field prediction and $37.5\%$ for $C_D$:
\begin{equation}
    \lambda_{\mathrm{tar}}^{\mathrm{field}}=0.5+0.5\frac{\mathbf{1}_{H_{0.1}}}{\rho(H_{0.1})},
    \qquad
    \lambda_{\mathrm{tar}}^{C_D}=0.625+0.375\frac{\mathbf{1}_{H_{0.1}}}{\rho(H_{0.1})}.
    \label{eq:app_pi_car_density}
\end{equation}

\paragraph{DrivAerML.}
Normal variation identifies geometric detail relevant to pressure reconstruction, while projected surface area measures sensitivity to pressure drag. We compute $s_n(x_i)=16^{-1}\sum_{j\in\mathcal{N}_{16}(i)}\|\mathbf{n}_i-\mathbf{n}_j\|_2^2$ and $s_D(x_i)=A_i|n_{i,x}|$, then increase density with the normalized indicator:
\begin{equation}
    \lambda_{\mathrm{tar}}^{\mathrm{field}}=\mathcal{N}[1+40\mathcal{T}_{0,0.99}[s_n]],
    \qquad
    \lambda_{\mathrm{tar}}^{C_D}=\mathcal{N}[1+40\mathcal{T}_{0,0.99}[s_D]].
    \label{eq:app_pi_drivaerml_density}
\end{equation}

\paragraph{SuperWing.}
Changes in surface normals identify curved or sharp wing features, motivating denser sampling around them. We compute $s_n=(\|\partial_I\mathbf{n}\|_2^2+\|\partial_J\mathbf{n}\|_2^2)^{1/2}$ on the surface lattice and normalize it as $\overline s_n=\mathcal{T}_{0.01,0.99}[s_n]$. The metric-specific densities are
\begin{equation}
    \lambda_{\mathrm{tar}}^{\mathrm{field}}=\mathcal{N}[1+16\overline s_n^2],
    \qquad
    \lambda_{\mathrm{tar}}^{C_D/C_L}=\mathcal{N}[1+8\overline s_n^{1/2}].
    \label{eq:app_pi_superwing_density}
\end{equation}

Plasticity, both Airfoil metrics, DrivAerML, and SuperWing use $\lambda_V$ compensation with $\beta=1$. Elasticity and Darcy use the target-density weights with $\beta=0.5$. For the two-stage samplers, let $H_d^+$ include both the wall band and any supplementary points. Giving all these points the same near-wall weight yields the equivalent compensation density $\lambda_{\mathrm{shell}}=\mathcal{N}[(1-\gamma)+\gamma\mathbf{1}_{H_d^+}/\rho(H_d)]$. Its inverse is raised to $\beta=0.75$ for Pipe and $\beta=0.5$ for ShapeNet-Car.

\subsection{Coarse-Prediction-Guided Density Refinement}
\label{app:cpg_density_design}

CPG evaluates the frozen operator on a uniform coarse set $S_0$ containing approximately $10\%$ of the reference points. The draft $\widehat u^{(0)}$ then guides density allocation. Unless noted otherwise, $S_0$ is retained and compensation uses $\lambda_V$ with $\beta=1$. Nearest-point cells use reference-lattice coordinates when available, and physical coordinates otherwise.

\paragraph{Elasticity.}
Stress variations identify regions where the response changes rapidly. We compute $s=\mathcal{I}_8D_8[\widehat\sigma_{\mathrm{vM}}^{(0)}]$ in physical coordinates and clip its upper tail. The target is realized through proportional quotas in 32 regions, with density contrast $\kappa=4$ and farthest-point placement:
\begin{equation}
    \lambda_{\mathrm{tar}}
    =0.375+0.625\mathcal{N}[\min(s,Q_{0.9}(s))].
    \label{eq:app_cpg_elasticity_density}
\end{equation}

\paragraph{Plasticity.}
Displacement gradients and curvature locate concentrated deformation throughout the loading trajectory. A quadratic fit over 16 neighbors gives $G=\|\nabla\widehat{\mathbf u}^{(0)}\|_F$ and $L=\|\Delta\widehat{\mathbf u}^{(0)}\|_2$ across all 40 displacement outputs, using physical spacing $0.5$. With coarse-point means $\overline G_0$ and $\overline L_0$, we transfer their normalized mixture and average over the full reference points in each region:
\begin{equation}
    s=\mathcal{I}_0\!\left[\frac{G}{2\overline G_0}+\frac{L}{2\overline L_0}\right],
    \qquad
    \lambda_{\mathrm{tar}}=\mathcal{P}_{24,4}[0.2+0.8\mathcal{N}[s]].
    \label{eq:app_cpg_plasticity_density}
\end{equation}

\paragraph{Airfoil.}
Pressure gradients locate flow variation, while surface-pressure contributions determine the drag-to-lift ratio. For field prediction, use $s_f=\mathcal{I}_0G_{16}[\widehat p^{(0)}]$. For the ratio, let $\widehat D=\sum_\Gamma b_D\widehat p^{(0)}$ and $\widehat L=\sum_\Gamma b_L\widehat p^{(0)}$ be the unnormalized force sums from \eqref{eq:app_airfoil_coefficients}. Define ratio sensitivity $v_R$ and balanced lift contributions $v_B$ by
\begin{equation}
    v_R=\frac{|\widehat Lb_D-\widehat Db_L|}{\max(|\widehat L|,0.02)^2},
    \qquad
    v_B=\mathcal{Q}_{0.95}[(b_L\widehat p^{(0)})_+]
          +\mathcal{Q}_{0.95}[(-b_L\widehat p^{(0)})_+].
    \label{eq:app_density_airfoil_signals}
\end{equation}
Here $(\cdot)_+$ is the positive part and $\chi=\min(1,|\widehat L|/0.02)$ downweights unreliable ratio estimates. We retain all 121 surface nodes and resample the interior using
\begin{equation}
    \begin{aligned}
    \lambda_{\mathrm{tar}}^{\mathrm{field}}
       &=\mathcal{N}[(0.1+\mathcal{Q}_{0.95}[s_f])^{2}],\\
    \lambda_{\mathrm{tar}}^{C_D/C_L}
       &=\mathcal{N}\!\left[
          \bigl(0.1+\chi\mathcal{K}_{0.1}[v_R]
          +\mathcal{K}_{0.1}[|b_L|]+2\mathcal{K}_{0.1}[v_B]
          +0.5e^{-d_\Gamma/0.1}\bigr)^{1/2}\right].
    \end{aligned}
    \label{eq:app_density_airfoil_ratio}
\end{equation}

\paragraph{Pipe.}
Local velocity variation identifies flow regions requiring higher resolution. We compute gradient energy from the coarse horizontal velocity in reference-lattice coordinates, then average it over coarse points in each region:
\begin{equation}
    \lambda_{\mathrm{tar}}=\mathcal{P}_{24,4}[D_8[\widehat v_x^{(0)}]].
    \label{eq:app_cpg_pipe_density}
\end{equation}

\paragraph{Darcy.}
Pressure gradients and curvature reveal the response to heterogeneous permeability. On the normalized draft pressure, we combine coarse gradient energy with the lattice Laplacian of its nearest-neighbor extension:
\begin{equation}
    s_{\mathrm{Darcy}}=D_8[\widehat p^{(0)}]
       +|\Delta_{\mathrm{FD}}\mathcal{I}_0\widehat p^{(0)}|_{S_0},
    \qquad
    \lambda_{\mathrm{tar}}=\mathcal{P}_{32,2}[s_{\mathrm{Darcy}}].
    \label{eq:app_cpg_darcy_density}
\end{equation}

\paragraph{ShapeNet-Car.}
Pressure variation guides field resolution, and positive pressure-force contributions identify regions relevant to drag. We compute surface gradient energy on coarse points and extend it to the surrounding volume. For drag, $\widehat p_\Gamma^{(0)}$ is the pressure transferred to the full surface, $\overline p_0$ is the mean coarse surface pressure, and $b_{D,i}=-\sum_{f\ni i}A_fn_{f,z}/4$ follows \eqref{eq:app_car_drag}. The two scores are
\begin{equation}
    s_f=\mathcal{E}_{0.1}[D_4[\widehat p^{(0)}]],
    \qquad
    s_D=\mathcal{E}_{0.1}[(b_D(\widehat p_\Gamma^{(0)}-\overline p_0))_+].
    \label{eq:app_cpg_car_scores}
\end{equation}
For either score, set $s'=\max(\mathcal{Q}_{0.95}[s],0.1)$ and $z_s=\operatorname{clip}_{[0,1]}(s'/Q_{0.95}(s'|_\Gamma))$. We retain the surface and fill capped proportional volume quotas by farthest-point sampling, stratifying by wall distance and four score-quantile groups. Compensation uses $\lambda_P$, with $\beta=1$ for field prediction and $0.75$ for drag. For field and drag targets $q=f,D$, the density combines near-wall coverage with the pressure signal:
\begin{equation}
    \lambda_{\mathrm{tar}}^q
       =\mathcal{N}[(1+A_qe^{-d_\Gamma/0.1})(1+2z_{s_q})],
    \qquad
    (A_f,A_D)=(64,32),\quad q\in\{f,D\}.
    \label{eq:app_density_car}
\end{equation}

\paragraph{DrivAerML.}
Local pressure variation guides surface resolution, and projected area indicates its relevance to drag. We transfer the coarse graph-Laplacian score $s=\mathcal{I}_0L_{16}[\widehat p^{(0)}]$ and use weighted sampling with
\begin{equation}
    \lambda_{\mathrm{tar}}^{\mathrm{field}}
       =\mathcal{N}[1+40\mathcal{T}_{0,0.99}[s]],
    \qquad
    \lambda_{\mathrm{tar}}^{C_D}
       =\mathcal{N}[1+40\mathcal{T}_{0,0.99}[A|n_x|s]].
    \label{eq:app_cpg_drivaerml_density}
\end{equation}

\paragraph{SuperWing.}
Joint pressure and skin-friction curvature identifies localized aerodynamic structures. With $\mathbf u=(C_p,150C_{f,\tau},300C_{f,z})$, we compute $s=\mathcal{T}_{0.01,0.99}[\|\Delta_5\mathcal{I}\widehat{\mathbf u}^{(0)}\|_2]$ and use full-reference regional means. For $C_D/C_L$, we mix this signal with the normalized normal-variation prior $g=\overline s_n$:
\begin{equation}
    \begin{aligned}
    d_\omega&=(1-\omega)\mathcal{N}[1+8g^{1/2}]
                      +\omega\mathcal{N}[1+8s^{1/2}],\\
    \mathcal{B}_\omega[g,s]
       &=\mathcal{N}\!\left[\max\!\left(d_\omega,\frac{\max_xd_\omega(x)}{8}\right)\right].
    \end{aligned}
    \label{eq:app_density_wing_blend}
\end{equation}
The ratio policy uses weighted residual sampling. The two densities are
\begin{equation}
    \lambda_{\mathrm{tar}}^{\mathrm{field}}=\mathcal{P}_{64,4}[s],
    \qquad
    \lambda_{\mathrm{tar}}^{C_D/C_L}=\mathcal{B}_{0.75}[g,s].
    \label{eq:app_cpg_superwing_density}
\end{equation}

\subsection{Cache-Aware Learned Refinement}
\label{app:ca_density_design}

CA obtains one scalar per coarse point from a head $H_\phi(h^{(1)})$ on frozen first-block features. Field targets are standardized, while nonnegative score targets are standardized after a $\log(1+s)$ transform. All quantities below denote decoded predictions. MLP heads use GELU between linear layers. Unless stated otherwise, the coarse budget is $10\%$, the final set retains $S_0$, and compensation follows the corresponding CPG policy.

\paragraph{Elasticity.}
Large local stress variations identify regions requiring refinement. A $128\to4\to1$ head predicts stress, from which we compute $s=\mathcal{I}_8D_8[\widehat\sigma_{\mathrm{vM}}^{\mathrm{head}}]$. We take regional maxima $v_k=\max_{x\in P_k}s(x)$ and assign proportional quotas across 32 regions with $\kappa=2$:
\begin{equation}
    \lambda_{\mathrm{tar}}(x)
       =0.75+0.25\frac{\min(v_k,Q_{0.99}(v))^2}
                         {\sum_j\mu_j\min(v_j,Q_{0.99}(v))^2}
    \quad(x\in P_k).
    \label{eq:app_density_elasticity_ca}
\end{equation}

\paragraph{Plasticity.}
The displacement-gradient magnitude summarizes concentrated deformation over the full loading trajectory. A $128\to4\to32\to1$ head learns $G_{\mathrm{GT}}=(\sum_{t,c}\|\nabla u_{t,c}\|_2^2)^{1/2}$ from the training fields, using reference spacing $0.5$. We transfer its estimate and draw weighted additions with
\begin{equation}
    \lambda_{\mathrm{tar}}=0.2+0.8\mathcal{N}[\mathcal{I}_0\widehat G^{\mathrm{head}}].
    \label{eq:app_ca_plasticity_density}
\end{equation}

\paragraph{Airfoil.}
Pressure gradients identify flow variation relevant to field reconstruction and aerodynamic prediction. A $128\to8\to32\to1$ head predicts pressure, with density parameters selected for each metric. The coarse budget is $10\%$ for the field and $5\%$ for $C_D/C_L$, always including the 121 surface nodes. Denoting the pressure estimate for target $q$ by $\widehat p^{\mathrm{head},q}$ gives
\begin{equation}
    \begin{aligned}
    s_q&=\mathcal{I}_0G_{16}[\widehat p^{\mathrm{head},q}],
       \quad q\in\{\mathrm{field},C_D/C_L\},\\
    \lambda_{\mathrm{tar}}^{\mathrm{field}}
       &=\mathcal{N}[(0.01+\mathcal{Q}_{0.95}[s_{\mathrm{field}}])^2],\\
    \lambda_{\mathrm{tar}}^{C_D/C_L}
       &=\mathcal{N}[(0.1+\mathcal{Q}_{0.95}[s_{C_D/C_L}])^{1/2}].
    \end{aligned}
    \label{eq:app_ca_airfoil_density}
\end{equation}
The CPG and CA field-policy parameters are screened on 24 geometries with one sampling seed and selected on 48 geometries with three seeds under the original $C_L>0.01$ inclusion rule. The selected policies are frozen before evaluation on 48 additional confirmation geometries. They are then reused without retuning for the 110 test cases satisfying $C_L>0.02$, with sampling seeds 0, 1, and 2. Both field policies use $\beta=1$.

\paragraph{Pipe.}
The velocity-gradient magnitude provides a direct measure of local flow variation. A head with two Transolver blocks of width 128 learns $G=\|\nabla v_x\|_2$ from the training fields. Squaring the decoded estimate gives the allocation score:
\begin{equation}
    \lambda_{\mathrm{tar}}
       =\mathcal{P}_{24,4}[(\widehat G^{\mathrm{head}})^2].
    \label{eq:app_ca_pipe_density}
\end{equation}

\paragraph{Darcy.}
Pressure-gradient and curvature signals capture the spatially varying response to permeability. A head with two Transolver blocks of width 128 learns the frozen teacher score $s_{\mathrm{Darcy}}$ in \eqref{eq:app_cpg_darcy_density}, giving
\begin{equation}
    \lambda_{\mathrm{tar}}=\mathcal{P}_{32,2}[\widehat s_{\mathrm{Darcy}}^{\mathrm{head}}].
    \label{eq:app_ca_darcy_density}
\end{equation}

\paragraph{ShapeNet-Car.}
Surface pressure supplies both local-variation and drag-contribution signals. A $256\to64\to1$ head predicts pressure, which replaces the coarse operator prediction in \eqref{eq:app_cpg_car_scores}. The surface retention, volume quotas, and compensation remain those of CPG. With $z_s$ defined as above, the learned densities are
\begin{equation}
    \lambda_{\mathrm{tar}}^q
       =\mathcal{N}[(1+A_qe^{-d_\Gamma/0.1})(1+2z_{s_q^{\mathrm{head}}})],
    \qquad
    (A_f,A_D)=(64,32),\quad q\in\{f,D\}.
    \label{eq:app_ca_car_density}
\end{equation}

\paragraph{DrivAerML.}
The local pressure residual estimates where surface predictions vary most, and projected area determines its relevance to drag. A $256\to8\to32\to1$ head learns $L_{16}[\widehat p^{(0)}]$ from the frozen teacher. After transfer, $s=\mathcal{I}_0\widehat L^{\mathrm{head}}$ gives
\begin{equation}
    \lambda_{\mathrm{tar}}^{\mathrm{field}}
       =\mathcal{N}[1+512\mathcal{T}_{0,0.99}[s]^2],
    \qquad
    \lambda_{\mathrm{tar}}^{C_D}
       =\mathcal{N}[1+40\mathcal{T}_{0,0.99}[A|n_x|s]].
    \label{eq:app_ca_drivaerml_density}
\end{equation}

\paragraph{SuperWing.}
Joint pressure and skin-friction curvature provides a compact signal for aerodynamic refinement. A $128\to4\to1$ head learns the normalized score $\mathcal{T}_{0.01,0.99}[\|\Delta_5\mathbf u_{\mathrm{GT}}\|_2]$. With $s=\mathcal{T}_{0.01,0.99}[\mathcal{I}\widehat s^{\mathrm{head}}]$ and the same geometric prior $g$, field prediction uses full-reference regional averages, while the ratio uses residual systematic sampling in Morton order:
\begin{equation}
    \lambda_{\mathrm{tar}}^{\mathrm{field}}=\mathcal{P}_{64,4}[s^{1/2}],
    \qquad
    \lambda_{\mathrm{tar}}^{C_D/C_L}=\mathcal{B}_{0.5}[g,s].
    \label{eq:app_ca_superwing_density}
\end{equation}
